% ============================================================
% M2 轮4 成卷·S6 件3 —— 滚动测评卷（B）（16 题＝命制-滚B 全卷题面入卷）
% 2026-09-12｜执行：S6 成卷代理（Kimi K2.8 第三臂首派，看板九十四）
% 样式：卷式宏承 v11 测评卷 qp-* 六宏（原样拷贝、零改动）＋练习线题块；与滚A 同风格（排版照同风格拍板）。
% 题源：命制/命制-滚B-16题.md 全卷 16 题（复核-滚B-glm 16/16 键合、0 键离、0 不通过；与滚A 逐题错开已核）。
% 卷式（照命制件头注）：单选 6×5＝30｜多选 3×5＝15｜填空 3×5＝15｜解答 8＋10＋10＋12＝40；40 分钟 100 分。
% 全卷无图（三维图禁绘；滚B-2 为平面二维图，题面数据自足径省；滚B-15/16 题面已内嵌完整指令，文字自足）。
% 复核临界登记随卷落实：滚B-15(3)↔简7-5(2) 同族警示（跨件型不并卷，天然满足）；滚B 自查「余弦」误记已更正为正弦。
% 答案制A＝题答分离，卷末答案速查表。
% ============================================================
\documentclass[fontset=none]{ctexart}
\usepackage[pure]{qp-m3}
\ansblockgrayfalse   % 换装（测评/滚动＝8 开三栏件）：答案块一律括线模（方案草案§1.2.5/§5.1）
\providecommand{\ov}[1]{\overrightarrow{#1}}   % 答案册 body 局部宏随值串迁入（丁区 10 键值含 \ov）

% ------------------------------------------------------------
% 件型层 A｜页面：8 开横放三栏（承测评卷件型层）
% ------------------------------------------------------------
\geometry{paperwidth=399.8mm,paperheight=284.2mm,left=8.9mm,right=8.9mm,
  top=12.0mm,bottom=17.6mm,twoside,headheight=0pt,headsep=0pt,footskip=7.1mm}
\setlength{\topskip}{0pt}
\raggedbottom
\renewcommand{\normalsize}{\fontsize{10.5pt}{17.7pt}\selectfont}
\linespread{1}\normalsize
\setlength{\parindent}{0pt}
\setlength{\parskip}{0pt}
\newlength{\jpcolw}\setlength{\jpcolw}{120.0mm}
\newlength{\jpgap}\setlength{\jpgap}{11.0mm}
\xeCJKsetup{CJKglue={\hskip 0.04em plus 0.04em minus 0.01em}}

\definecolor{silklight}{gray}{0.8235}
\definecolor{silkdark}{gray}{0.6510}
\definecolor{juanrule}{gray}{0.1255}
\definecolor{subgrey}{gray}{0.4}

% ------------------------------------------------------------
% 件型层 B｜JT-01 丝带角饰（照抄测评卷件型层）
% ------------------------------------------------------------
\newcommand{\juanmathSize}{17pt}
\newcommand{\sishou}{%
  \begin{tikzpicture}[overlay,remember picture,x=1mm,y=1mm]
    \path (current page.north west) coordinate (O);
    \begin{scope}[shift={(O)}]
      \fill[silklight] (6.9,-9.39) -- (40.1,-9.39) -- (6.9,-42.59) -- cycle;
      \draw[white,line width=0.28mm] (9.21,-40.29) -- (9.21,-11.63) -- (37.87,-11.63);
      \fill[silkdark] (22.94,-8.21) -- (34.15,-8.21) -- (5.78,-36.58) -- (5.78,-25.37) -- cycle;
      \fill[silkdark] (34.15,-8.21) -- (34.15,-9.20) -- (33.16,-9.20) -- cycle;
      \fill[silkdark] (5.78,-36.58) -- (6.78,-36.58) -- (6.78,-35.58) -- cycle;
      \node[white,inner sep=0pt,anchor=center,rotate=45] at (17.23,-19.74)
        {\fontsize{\juanmathSize}{\juanmathSize}\selectfont\heitiao 数学};
    \end{scope}
  \end{tikzpicture}}

% ------------------------------------------------------------
% 件型层 C｜卷头零件＋题区零件（承测评卷件型层；卷名/副题/时间/组行按滚B 卷式改文）
% ------------------------------------------------------------
\newcommand{\zeroheight}[1]{\raisebox{0pt}[0pt][0pt]{#1}}
\newcommand{\jpcen}[1]{\hbox to\jpcolw{\hfil\zeroheight{#1}\hfil}\nointerlineskip}
\newcommand{\jpright}[1]{\hbox to\jpcolw{\hfill\zeroheight{#1}}\nointerlineskip}
\newcommand{\vgt}[1]{\par\vskip\dimexpr#1-6.25mm\relax}

\newcommand{\tianxieg}[1]{{\fontsize{\qpJTbLabelSize}{13pt}\selectfont\heitiao #1：}%
  \rule[-0.62mm]{\qpJTbLineLen}{0.2mm}}
\newcommand{\jtianxie}{\tianxieg{班级}\hspace{\qpJTbGapGrp}\tianxieg{姓名}%
  \hspace{\qpJTbGapGrp}\tianxieg{得分}}
\newcommand{\juanmingSize}{20.5pt}
\newcommand{\jjuanming}{{\fontsize{\juanmingSize}{26pt}\selectfont\heizhang
  \renewcommand{\CJKglue}{\hskip 0pt}滚动测评卷（B）}}
\newcommand{\jjunrule}{{\color{juanrule}\rule{58.0mm}{\qpJTcRule}}}
\newcommand{\jfuti}{{\fontsize{\qpJTdSize}{18pt}\selectfont\heitiao\color{subgrey}%
  \renewcommand{\CJKglue}{\hskip 0pt}第一章（全章）}}
\newcommand{\jptime}{{\fontsize{\qpJTeSize}{17.7pt}\selectfont（时间：40分钟\quad 分值：100分）}}
\newcommand{\zuhang}[2]{\par\noindent\hangindent=5.7mm\hangafter=1
  {\fontsize{\qpJTfSize}{17.7pt}\selectfont\heitiao #1}#2\par}
\newcommand{\ti}[2]{\par\noindent\hangindent=5.7mm\hangafter=1
  {\fontsize{11.4pt}{17.7pt}\selectfont\heihao{\numboldjian #1}.}\hspace{2.7mm}#2\par}
\newcommand{\fenzhi}[1]{（#1分）}
\newcommand{\optline}[1]{\par\noindent\hangindent=5.7mm\hangafter=0
  {\fontsize{10.5pt}{19pt}\selectfont #1}\par}
\newcommand{\optII}[2]{\makebox[53.05mm][l]{#1}#2}
\newcommand{\optIV}[4]{\makebox[21.5mm][l]{#1}\makebox[21.5mm][l]{#2}\makebox[21.5mm][l]{#3}#4}
\newcommand{\kw}{\hfill（\hspace{3.6mm}）\hspace*{5.4mm}}
\newcommand{\qpart}[1]{\par\noindent\hangindent=5.7mm\hangafter=0
  \hspace*{5.7mm}{\fontsize{10.5pt}{17.7pt}\selectfont #1}\par}

\newdimen\hA \hA=3.03mm
\newdimen\hB \hB=12.12mm
\newdimen\hC \hC=1.89mm
\newdimen\hD \hD=5.95mm
\newdimen\hE \hE=8.04mm
\newdimen\hF \hF=5.70mm
\newdimen\hG \hG=6.35mm
\newcommand{\jphead}{\hbox to\jpcolw{\sishou\hfill}\nointerlineskip
  \vskip\hA\jpright{\jtianxie}%
  \vskip\hB\jpcen{\jjuanming}%
  \vskip\hC\jpcen{\jjunrule}%
  \vskip\hD\jpcen{\jfuti}%
  \vskip\hE\jpcen{\jptime}%
  \vskip\hF
  \zuhang{一、选择题}{：本题共6小题，每小题5分，共30分．\ 在每小题给出的四个选项中，只有一项是符合题目要求的．}%
  \vgt{\hG}}

% ------------------------------------------------------------
% 件型层 D｜YJ-02 文字行页脚（同测评卷；卷名改滚B）
% ------------------------------------------------------------
\newcommand{\jpnum}[1]{{\fontsize{\qpYJbNumSize}{\qpYJbNumSize}\selectfont\numbold #1}}
\newcommand{\jptxtsize}{7.5pt}
\newcommand{\jptxt}{\fontsize{\jptxtsize}{\jptxtsize}\selectfont}
\newcommand{\juanname}{滚动测评卷（B）}
\newcommand{\pqt}{\ifnum\value{page}<10 0\fi\thepage}
\setlength{\headwidth}{\textwidth}
\fancyfoot[RO]{\jptxt\juanname\hspace{3.95mm}{\heitiao 滚动卷}\hspace{3.88mm}卷\hspace{2.25mm}\jpnum{\pqt}}
\fancyfoot[LE]{\jptxt\jpnum{\pqt}\hspace{1.9mm}卷\hspace{3.95mm}{\heitiao 羿郭工作室}%
  \hspace{3.0mm}高中数学\hspace{2.75mm}选择性必修第一册\hspace{2.8mm}RJB}

% ------------------------------------------------------------
% 件型层 E｜三栏排布器＋卷末答案速查表
% ------------------------------------------------------------
\newcommand{\jpcol}[1]{\raisebox{0pt}[0pt][0pt]{\vtop{\hsize\jpcolw\parskip0pt\parindent0pt
  \setlength\linewidth{\jpcolw}\setlength\columnwidth{\jpcolw}   % 换装补钉：qp-m3 括线盒按 \linewidth 取宽
  \fontsize{10.5pt}{17.7pt}\selectfont\relax#1}}}
\newcommand{\jpthree}[3]{\nointerlineskip\noindent\jpcol{#1}\hspace{\jpgap}\jpcol{#2}%
  \hspace{\jpgap}\jpcol{#3}\par}
\newcommand{\jplead}{\hbox{}\nointerlineskip}

\newcommand{\datu}[1]{{\fontsize{9pt}{12pt}\selectfont #1}}
\newcommand{\dabiao}{%
  \par\vgt{\hG}\noindent{\fontsize{10.5pt}{17.7pt}\selectfont\heitiao 答案速查}\par\nopagebreak
  \vskip 1.2mm\noindent\datu{\setlength{\tabcolsep}{1.4pt}%
  \begin{tabular}{|c|*{16}{c|}}
  \hline
  题号 & 1 & 2 & 3 & 4 & 5 & 6 & 7 & 8 & 9 & 10 & 11 & 12 & 13 & 14 & 15 & 16 \\ \hline
  答案 & A & A & B & A & C & C & AC & ABC & AB & \(0\) & \(\frac{7}{2}\) & \(\frac{2\sqrt{3}}{3}\) & 略 & 略 & 略 & 略 \\ \hline
  \end{tabular}}\par}

\begin{document}
% ===================== 第 1 页：栏1 卷头＋Q1–Q3；栏2 Q4–Q6；栏3 组行二＋Q7–Q9 =====================
\jpthree{\jphead
  \ti{1}{已知向量 \(\overrightarrow{a}\)，\(\overrightarrow{b}\) 满足 \(\left|\overrightarrow{a}\right|=2\)，\(\left|\overrightarrow{b}\right|=1\)，\(\langle\overrightarrow{a},\overrightarrow{b}\rangle=120^\circ\)，则 \(\overrightarrow{a}\cdot(\overrightarrow{a}+2\overrightarrow{b})=\)\kw}
  \optline{\optIV{A．\(2\)；}{B．\(-2\)；}{C．\(6\)；}{D．\(8\)．}}
  \ti{2}{在 \(\triangle ABC\) 中，\(\angle A=90^\circ\)，\(AB=2\)，\(AC=1\)，则 \(\overrightarrow{BC}\cdot\overrightarrow{AB}=\)\kw}
  \optline{\optIV{A．\(-4\)；}{B．\(-2\)；}{C．\(2\)；}{D．\(4\)．}}
  \ti{3}{设 \(\{\overrightarrow{e_1},\overrightarrow{e_2},\overrightarrow{e_3}\}\) 是空间的一个基底，\(\overrightarrow{a}=\overrightarrow{e_1}+\overrightarrow{e_2}\)，\(\overrightarrow{b}=\overrightarrow{e_2}+\overrightarrow{e_3}\)，则下列向量中能与 \(\overrightarrow{a}\)，\(\overrightarrow{b}\) 构成空间另一个基底的是\kw}
  \optline{\optII{A．\(\overrightarrow{e_1}-\overrightarrow{e_3}\)；}{B．\(\overrightarrow{e_1}+\overrightarrow{e_2}+\overrightarrow{e_3}\)；}}
  \optline{\optII{C．\(\overrightarrow{a}+2\overrightarrow{b}\)；}{D．\(\overrightarrow{a}+\overrightarrow{b}\)．}}
}{\jplead
  \ti{4}{设空间中一点 \(S(1,2,-1)\)，下列各点中，到点 \(S\) 的距离等于 \(3\) 的是\kw}
  \optline{\optIV{A．\((1,2,2)\)；}{B．\((0,0,0)\)；}{C．\((1,2,-2)\)；}{D．\((2,1,1)\)．}}
  \ti{5}{直线 \(l\) 与平面 \(\alpha\) 相交于点 \(B\)，点 \(A\notin\alpha\)，\(AA'\perp\alpha\) 于点 \(A'\)．若 \(AA'=2\sqrt{3}\)，\(A'B=2\)，则直线 \(AB\) 与平面 \(\alpha\) 所成角的大小为\kw}
  \optline{\optIV{A．\(30^\circ\)；}{B．\(45^\circ\)；}{C．\(60^\circ\)；}{D．\(90^\circ\)．}}
  \ti{6}{已知平面 \(\alpha\) 与平面 \(\beta\) 的法向量的夹角为 \(60^\circ\)，且二面角 \(\alpha-l-\beta\) 为钝角，则该二面角的大小是\kw}
  \optline{\optIV{A．\(30^\circ\)；}{B．\(60^\circ\)；}{C．\(120^\circ\)；}{D．\(150^\circ\)．}}
}{\jplead
  \zuhang{二、选择题}{：本题共3小题，每小题5分，共15分．\ 在每小题给出的选项中，有多项符合题目要求．（全部选对得5分，部分选对得2分，有选错的得0分）}%
  \vgt{\hG}
  \ti{7}{下列说法正确的是\kw}
  \optline{A．两条平行直线的方向向量必定共线；}
  \optline{B．若空间中两条直线没有公共点，则这两条直线互相平行；}
  \optline{C．两条直线所成角的余弦值，等于它们方向向量夹角余弦值的绝对值；}
  \optline{D．若一条直线与一个平面内的无数条直线都垂直，则这条直线与该平面垂直．}
  \ti{8}{在棱长为 \(1\) 的正方体 \(ABCD-A_1B_1C_1D_1\) 中，下列结论正确的是\kw}
  \optline{A．异面直线 \(A_1B\) 与 \(B_1C\) 所成的角为 \(60^\circ\)；}
  \optline{B．\(AC_1\perp BD\)；}
  \optline{C．直线 \(AA_1\) 与 \(BC_1\) 所成的角为 \(45^\circ\)；}
  \optline{D．向量 \(\overrightarrow{AB_1}\) 的模等于 \(1\)．}
  \ti{9}{在棱长为 \(1\) 的正方体 \(ABCD-A_1B_1C_1D_1\) 中，下列有关二面角的结论正确的是\kw}
  \optline{A．二面角 \(D-AB-A_1\) 是直二面角；}
  \optline{B．二面角 \(C_1-BD-C\) 的余弦值为 \(\frac{\sqrt{3}}{3}\)；}
  \optline{C．二面角 \(B_1-AC-D_1\) 的余弦值为 \(\frac{\sqrt{3}}{3}\)；}
  \optline{D．二面角 \(A_1-CD-B\) 的大小为 \(60^\circ\)．}
}
\newpage
% ===================== 第 2 页：栏1 组行三＋Q10–Q12＋组行四；栏2 Q13–Q14；栏3 Q15–Q16＋答案速查表 =====================
\jpthree{\jplead
  \zuhang{三、填空题}{：本题共3小题，每小题5分，共15分．}%
  \vgt{\hG}
  \ti{10}{已知向量 \(\overrightarrow{a}=(1,2,2)\)，\(\overrightarrow{b}=(1,-2,2)\)，则向量 \(\overrightarrow{a}+\overrightarrow{b}\) 与 \(\overrightarrow{a}-\overrightarrow{b}\) 夹角的余弦值为\kongbai}
  \ti{11}{平面直角坐标系中，三点 \(A(1,2)\)，\(B(5,10)\)，\(C(3,2k-1)\) 在同一条直线上，则实数 \(k=\)\kongbai}
  \ti{12}{棱长为 \(2\) 的正方体 \(ABCD-A_1B_1C_1D_1\) 中，点 \(E\) 在棱 \(DD_1\) 上运动，则点 \(E\) 到平面 \(A_1BD\) 的距离与点 \(E\) 到平面 \(AB_1D_1\) 的距离之和等于\kongbai}
  \vgt{\hG}
  \zuhang{四、解答题}{：本题共4小题，共40分．\ 解答应写出文字说明、证明过程或演算步骤．}%
  \vgt{\hG}
}{\jplead
  \ti{13}{\fenzhi{8}长方体 \(ABCD-A_1B_1C_1D_1\) 中，\(AB=AD=2\)，\(AA_1=4\)，\(E\) 为棱 \(BB_1\) 的中点．以 \(D\) 为原点，射线 \(DA\)，\(DC\)，\(DD_1\) 分别为 \(x\) 轴、\(y\) 轴、\(z\) 轴的正方向建立空间直角坐标系．}
  \qpart{（1）写出点 \(A\)，\(B_1\)，\(E\) 的坐标；}
  \qpart{（2）求向量 \(\overrightarrow{AE}\) 与 \(\overrightarrow{CB_1}\) 夹角的余弦值．}
  \ti{14}{\fenzhi{10}已知空间中三点 \(A(1,0,0)\)，\(B(0,2,0)\)，\(C(0,0,3)\)．}
  \qpart{（1）求平面 \(ABC\) 的一个法向量；}
  \qpart{（2）若点 \(Q(1,t,3)\)（\(t\in\mathbf{R}\)）在平面 \(ABC\) 内，求 \(t\) 的值；}
  \qpart{（3）求点 \(D(2,1,2)\) 到平面 \(ABC\) 的距离．}
}{\jplead
  \ti{15}{\fenzhi{10}三棱锥 \(P-ABC\) 中，\(PA\perp\) 平面 \(ABC\)，\(AB\perp BC\)，\(PA=2\)，\(BC=2\)．}
  \qpart{（1）求证：\(BC\perp\) 平面 \(PAB\)；并在平面 \(PAB\) 内过 \(A\) 作 \(AE\perp PB\) 于 \(E\)，说明线段 \(AE\) 的长就是点 \(A\) 到平面 \(PBC\) 的距离；}
  \qpart{（2）若 \(AE=\sqrt{3}\)，求 \(AB\) 的长；}
  \qpart{（3）在（2）的条件下，求直线 \(PC\) 与平面 \(PAB\) 所成角的正弦值．}
  \ti{16}{\fenzhi{12}四面体 \(ABCD\) 中，\(AB=BC=AD=DC=2\)，\(AC=2\sqrt{3}\)，\(BD=\sqrt{3}\)．}
  \qpart{（1）求二面角 \(B-AC-D\) 的大小；}
  \qpart{（2）求点 \(B\) 到平面 \(ACD\) 的距离；}
  \qpart{（3）求直线 \(AB\) 与平面 \(ACD\) 所成角的正弦值．}
  \ifshowans\dabiao\fi   % 换装案B：速查表＝答案承载件，纯题档随开关吞（防漏答案）
  \ifshowans\else % 案B双锚补丁：附卷页整页跳过，纯题档在此补偿发射 16 键（两档 M3-ANSKEY 恒等，对号门两档可跑）
  \anskey{滚B-1}\anskey{滚B-2}\anskey{滚B-3}\anskey{滚B-4}\anskey{滚B-5}\anskey{滚B-6}\anskey{滚B-7}\anskey{滚B-8}\anskey{滚B-9}\anskey{滚B-10}\anskey{滚B-11}\anskey{滚B-12}\anskey{滚B-13}\anskey{滚B-14}\anskey{滚B-15}\anskey{滚B-16}
  \fi
}
% ===================== 卷末附卷：参考答案（案B 附卷制：仅含详解档成页；纯题档整页跳过） =====================
\ifshowans
\newpage
\jpthree{\jplead
  \par\vskip 2.0mm\noindent{\fontsize{\qpJTfSize}{17.7pt}\selectfont\heitiao 参考答案}\hspace{0.6em}{\fontsize{\qpJTeSize}{17.7pt}\selectfont（滚动测评卷（B）·含详解档）}\par\nopagebreak
  \begin{ansblock}[滚B-1]
  % ans:滚B-1
  \ansitem{1}{A}
  \end{ansblock}
  \begin{ansblock}[滚B-2]
  % ans:滚B-2
  \ansitem{2}{A}
  \end{ansblock}
  \begin{ansblock}[滚B-3]
  % ans:滚B-3
  \ansitem{3}{B}
  \end{ansblock}
  \begin{ansblock}[滚B-4]
  % ans:滚B-4
  \ansitem{4}{A}
  \end{ansblock}
  \begin{ansblock}[滚B-5]
  % ans:滚B-5
  \ansitem{5}{C}
  \end{ansblock}
  \begin{ansblock}[滚B-6]
  % ans:滚B-6
  \ansitem{6}{C}
  \end{ansblock}
}{\jplead
  \begin{ansblock}[滚B-7]
  % ans:滚B-7
  \ansitem{7}{AC}
  \end{ansblock}
  \begin{ansblock}[滚B-8]
  % ans:滚B-8
  \ansitem{8}{ABC}
  \end{ansblock}
  \begin{ansblock}[滚B-9]
  % ans:滚B-9
  \ansitem{9}{AB}
  \end{ansblock}
  \begin{ansblock}[滚B-10]
  % ans:滚B-10
  \ansitem{10}{\((0)\)}
  \end{ansblock}
  \begin{ansblock}[滚B-11]
  % ans:滚B-11
  \ansitem{11}{\(\dfrac{7}{2}\)}
  \end{ansblock}
  \begin{ansblock}[滚B-12]
  % ans:滚B-12
  \ansitem{12}{\(\dfrac{2\sqrt{3}}{3}\)}
  \end{ansblock}
}{\jplead
  \begin{ansblock}[滚B-13]
  % ans:滚B-13
  \ansitem{13}{(1)\(A(2,0,0)\)、\(B_{1}(2,2,4)\)、\(E(2,2,2)\)；(2)\(\dfrac{\sqrt{10}}{5}\)}
  \ansnote{解析}{(2) \(\ov{AE}=(0,2,2)\)、\(\ov{CB_{1}}=(2,0,4)\)，\(\cos\langle\ov{AE},\ov{CB_{1}}\rangle=\dfrac{8}{2\sqrt{2}\cdot 2\sqrt{5}}=\dfrac{\sqrt{10}}{5}\)。}
  \end{ansblock}
  \begin{ansblock}[滚B-14]
  % ans:滚B-14
  \ansitem{14}{(1)\(\ov{n}=(6,3,2)\)；(2)\(t=-2\)；(3)\(\dfrac{13}{7}\)}
  \ansnote{解析}{(1) \(\ov{AB}=(-1,2,0)\)、\(\ov{AC}=(-1,0,3)\)，联立 \(\ov{n}\cdot\ov{AB}=\ov{n}\cdot\ov{AC}=0\) 取 \(\ov{n}=(6,3,2)\)。(2) \(\ov{AQ}=(0,t,3)\)，\(\ov{n}\cdot\ov{AQ}=3t+6=0\)，\(t=-2\)。(3) \(\ov{AD}=(1,1,2)\)，\(d=\dfrac{|\ov{n}\cdot\ov{AD}|}{|\ov{n}|}=\dfrac{13}{\sqrt{49}}=\dfrac{13}{7}\)。}
  \end{ansblock}
  \begin{ansblock}[滚B-15]
  % ans:滚B-15
  \ansitem{15}{(1)\(BC\perp\) 面 \(PAB\)、\(AE\perp\) 面 \(PBC\)（\(AE\) 长即距离）；(2)\(AB=2\sqrt{3}\)；(3)\(\dfrac{\sqrt{5}}{5}\)}
  \ansnote{解析}{(1) \(PA\perp\) 面 \(ABC\) 得 \(PA\perp BC\)，又 \(AB\perp BC\)，故 \(BC\perp\) 平面 \(PAB\)；于是 \(BC\perp AE\)，结合 \(AE\perp PB\) 得 \(AE\perp\) 平面 \(PBC\)，垂足在面内——\(AE\) 长即点 \(A\) 到面 \(PBC\) 的距离。(2) 设 \(AB=x\)：\(\mathrm{Rt}\triangle PAB\) 中 \(AE=\dfrac{PA\cdot AB}{PB}=\dfrac{2x}{\sqrt{x^{2}+4}}=\sqrt{3}\)，解得 \(x^{2}=12\)，\(AB=2\sqrt{3}\)。(3) \(PB\) 为 \(PC\) 在面 \(PAB\) 内的射影，\(\sin\angle CPB=\dfrac{BC}{PC}=\dfrac{2}{\sqrt{16+4}}=\dfrac{\sqrt{5}}{5}\)。}
  \end{ansblock}
  \begin{ansblock}[滚B-16]
  % ans:滚B-16
  \ansitem{16}{(1)\(120^{\circ}\)；(2)\(\dfrac{\sqrt{3}}{2}\)；(3)\(\dfrac{\sqrt{3}}{4}\)}
  \ansnote{解析}{(1) 取 \(AC\) 中点 \(O\)：\(BO\perp AC\)、\(DO\perp AC\)，\(BO=DO=1\)，\(\cos\angle BOD=\dfrac{1+1-3}{2\cdot 1\cdot 1}=-\dfrac{1}{2}\)，二面角 \(B\text{-}AC\text{-}D\) 为 \(120^{\circ}\)。(2) \(AC\perp\) 平面 \(BOD\)，面 \(BOD\perp\) 面 \(ACD\)、交线 \(OD\)；作 \(BH\perp OD\)，则 \(BH\perp\) 面 \(ACD\)，\(BH=BO\sin 60^{\circ}=\dfrac{\sqrt{3}}{2}\) 即距离。(3) \(\sin\theta=\dfrac{BH}{AB}=\dfrac{\sqrt{3}/2}{2}=\dfrac{\sqrt{3}}{4}\)。}
  \end{ansblock}
}
\fi

\end{document}
